\section{States, Mediators and Clock Comparisons in the 4.10 Research Line}
\label{sec:v410-state-mediator}

This section consolidates the investigations of 28--30 October 2026.
It distinguishes exact local identities, finite numerical evidence and
open physical identifications. The common carrier is a research programme;
the comparisons below neither close OP5 nor derive a unified quantum/gravity
action. The QAM records the types, hypotheses and remaining obligations of
an argument. An open entry is not a refutation, and a successful residual
test is not a substitute for the missing implication.

\subsection{A common carrier and distinct readouts}
\label{subsec:v410-readout-types}

Let \(\mathcal C\) be a declared common state carrier and let
\(\pi_Q:\mathcal C\to\mathcal Q\) and
\(\pi_R:\mathcal C\to\mathcal R\) be quantum and relativistic readouts
on specified admissible domains. A mediator \(\mathcal M\) requires
explicit maps \(m_Q,m_R\) with
\[
 m_Q\pi_Q(c)=m_R\pi_R(c).
\]
If dynamics are also claimed, the additional conditions are
\(\pi_Q U_{\mathcal C}=V_Q\pi_Q\) and
\(\pi_R U_{\mathcal C}=V_R\pi_R\), with compatible clock parameters.
These are proof obligations, not definitions that establish a physical
identification by themselves. A physical state space, a local null-direction
sphere, and a normalized two-port register retain different information.

In particular, the ordinary positive AB Gram register of two Dirac spinors
is not the same readout as the real electromagnetic doublet of r04.
Writing \(N=\|A\|^2+\|B\|^2>0\), the former gives
\[
 \rho_{AB}=\frac1N
 \begin{pmatrix}A^\dagger A&B^\dagger A\\A^\dagger B&B^\dagger B\end{pmatrix},
 \qquad
 r=\frac{(2\operatorname{Re}A^\dagger B,
                  2\operatorname{Im}A^\dagger B,
                  \|A\|^2-\|B\|^2)}N.
\]
It obeys \(1-|r|^2=4\det\rho_{AB}\ge0\). Proportional AB fields lie
on the boundary; independent fields can lie in the interior. Two fields
with the same \(r\) can have opposite
\(A^\dagger\beta A-B^\dagger\beta B\): take
\((A,B)=(\cos\theta\,e_1,\sin\theta\,e_1)\) or
\((\cos\theta\,e_3,\sin\theta\,e_3)\) with
\(\beta=\operatorname{diag}(1,1,-1,-1)\).
Thus the Bloch coordinates alone do not close the Dirac readout.

\subsection{An orientation-complete local mediator}
\label{subsec:v410-oriented-mediator}

For the native real doublet \(F=[E\mid B]\), put
\(G=F^{\mathsf T}F=\left(\begin{smallmatrix}a&c\\c&b\end{smallmatrix}\right)\).
On the chart \(a>0\), \(\Delta=\sqrt{ab-c^2}>0\), define
\[
 F_0(G)=\begin{pmatrix}\sqrt a&c/\sqrt a\\0&\Delta/\sqrt a\\0&0\end{pmatrix},
 \quad e=E/\sqrt a,\quad
 q=\frac{B-(c/a)E}{\Delta/\sqrt a},\quad O=[e\mid q\mid e\times q].
\]
Then \(O\in SO(3)\) and \(F=OF_0(G)\), by direct multiplication.
The three Gram entries and three orientation parameters therefore recover
the six real field components locally. Keeping only \(\rho=G/\operatorname{tr}G\)
also discards the scale \(\operatorname{tr}G=2u\). Rank-deficient fields
need other charts and stabilizer data. A nonzero null electromagnetic
field with \(G=uI\) is regular in this rank-two chart; it is not a
rank-one Gram state.

For any declared right-doublet operation \(M\),
\[
 F'=FM,\qquad G'=M^{\mathsf T}GM.
\]
Reconstructing \(F_0(G)M=K F_0(G')\) gives \(O'=OK\) whenever the
rank-two charts exist. This proves that changing the field and then
reading it agrees with the corresponding mediator operation. Exact checks
cover phase rotation, \(J_B=i\) times conjugation, conjugation and the
compensator \(C_+=(I+\sigma_x)/2\).

The compensator is a separate rank-changing case:
\(E'=B'=(E+B)/2\). If nonzero, its normalized Gram state is
\(\rho'=C_+\), and a rank-one direction chart replaces the orientation
triad. If the selected port vanishes, normalization is undefined.
The pure selected register state is not the same readout as an
electromagnetic null field, whose r04 state is \(I/2\).

Orientation is needed for general observer changes. For example,
\(E=(2,0,0)\), \(B_\pm=(0,\pm1,0)\) give the same
\(\rho=\operatorname{diag}(4/5,1/5)\), but opposite fluxes. Under a
common infinitesimal boost of speed \(v\) along \(z\),
\[
 \left.\frac{d\rho_+}{dv}\right|_0
 =\operatorname{diag}(12/25,-12/25)
 =-\left.\frac{d\rho_-}{dv}\right|_0.
\]
Even the flux alone does not repair every ambiguity: two orthogonal pure
electric fields have the same Gram matrix and zero flux but respond
differently to a boost along one of their axes. These facts establish
non-autonomy of the reduced register under unrestricted directed observer
changes, not a failure of the full field's covariance.

\subsection{Clock calibration and the octonionic scope}
\label{subsec:v410-clock-scope}

Coordinate time, eigen-time along a timelike observer, an internal clock
reading, and a structural invariant require separate definitions. In the
controlled constant-rate example
\(\psi_a(t_a)=\exp(-i\omega a t_a)\psi_0\),
\(\psi_b(t_b)=\exp(-i\omega b t_b)\psi_0\), \(a,b>0\), the calibration
\(\tau=a t_a=b t_b\) makes the trajectories agree exactly. Matching
their initial zeros alone does not. The rates are inputs of this example,
not derived native clocks. The orientation counterexample above persists
under a common calibrated clock whenever the observer's boost rate is
nonzero. Neither result requires a second fundamental quantum time.

A shared clock does not alone prove entanglement: local evolution
\(U_1(\tau)\otimes U_2(\tau)\) preserves product states. Likewise, an
invariant \(\tau_n\) is not by itself an operational clock or an
entanglement criterion. Quantum systems can also be massive; the two
descriptions are not disjoint massless/massive classes.

After choosing a complex direction, a six-real-dimensional octonionic
complement can encode \(E+iB\). This is a data embedding, not yet an
identification of Lorentz boosts with real octonion automorphisms.
Those preserve the positive octonion norm, whereas a transverse boost of
\(E=(0,1,0),B=0\) gives
\(E'^2-B'^2=1\) but
\(E'^2+B'^2=(1+v^2)/(1-v^2)\).
The native complex/quaternionic/octonionic dictionary and its allowed
actions must therefore be specified. Extra registers or complexification
are possible research routes, not consequences of dimension counting.

\begin{samepage}
\subsection{Field compatibility and finite evidence}
\label{subsec:v410-field-compatibility}

The declared classical comparison retains
\[
 \begin{aligned}
 F_\Psi&=i\sum_{\mu,a}(E^{-1})^\mu{}_a(I\otimes K_a)D_\mu\Psi
             -h(Y\otimes\beta)\Psi=0,\\
 E^a{}_{\mu}&=\operatorname{Re}(\Psi^\dagger Q_aD_\mu\Psi),\qquad\det E\ne0.
 \end{aligned}
\]
\end{samepage}
where \(Q_a=S\otimes K_a\), \(K_a=(I,\alpha_i)\), and the Lorentz
spin connection is set to zero as a comparison choice, not identified
with the induced Levi--Civita connection. The source couplings are held
fixed. The earlier proportional AB ansatz need not be preserved by this
dynamics. Releasing the first jets removes several pointwise residuals,
but does not establish their regional continuation.

For constant Hermitian \(Q_a\) and Hermitian internal connection,
\[
 E^a=\tfrac12dj_a+\tfrac i2\Psi^\dagger[Q_a,A]\Psi,
 \qquad j_a=\Psi^\dagger Q_a\Psi.
\]
For the retained \(A=G\rho\,d\varphi\) on \(\rho>0\), this implies
\(E^a=dj_a/2+B_a d\varphi\) and
\(dE^a=dB_a\wedge d\varphi\). Here the radial coordinate \(\rho\)
is unrelated to a density matrix. This coordinate curl is not itself
Riemann curvature. Eliminating all free ordered second jets shows
numerically that imposing even the full current-curl identity does not
reproduce the obstruction from fully symmetric spinor second jets.
Coframe integrability and spinor integrability are different tests.

At fixed field values let \(A_0\) be the real map from time-jet changes
to changes of \(E_0\), and \(B_0\) have the four covariant jets as real
columns. Then \(A_0B_0=E\). Pointwise shear invariance gives, for
invertible frozen time operator \(H_0\),
\[
 P=I-B_0E^{-1}A_0,\quad J_0=H_0P,\quad
 L=A_0H_0^{-1},\quad LJ_0=0.
\]
Thus \(LF_\Psi=0\) supplies four temporal compatibility conditions.
One combination is
\(4-h\Psi^\dagger(Z\otimes\beta)\Psi/2=0\),
\(Z=\operatorname{diag}(1,-1,0)\). Its derivatives and the other
compatibility combinations can be solved jointly at the three tested
points. The resulting first-jet family has a numerically regular
seven-direction constraint Jacobian and 65 free spatial tangent
directions, before further constraints and identifications. This is
neither a count of physical degrees of freedom nor a regional existence
theorem. For these frozen first jets, symmetric second-jet completion
retains nonzero residuals; the joint recalibration below releases those
data and supersedes this as a current obstruction.

Two-mode fields with variable amplitudes provide single smooth
approximations and hence consistent mixed derivatives, but retain
nonzero equation residuals. Their optimization is underdetermined and
does not supply a common action. Figure~\ref{fig:op5-multimode-slices}
deliberately records the original affine baseline, not the latest fit.
Later quadratic fits and weighted comparisons are archived separately.
The early edge-weighting interpretation was corrected when a reused
residual function was found to retain the old points; only the subsequent
corrected comparisons support statements about edge weighting.

As a precision control, the direct centre-point problem was recomputed
with necessary state normalization restored in each precision:
\begin{center}
\begin{tabular}{@{}rrr@{}}\toprule
decimal digits & normalized-input full residual & rounded-input full residual\\
\midrule
30 & \(2.3\times10^{-30}\) & \(2.9\times10^{-15}\)\\
50 & \(6.0\times10^{-50}\) & \(2.9\times10^{-15}\)\\
80 & \(4.5\times10^{-80}\) & \(2.9\times10^{-15}\)\\
\bottomrule\end{tabular}
\end{center}
These are numerical checks, not certified exact solutions or physical
scales. Renormalization changes the rounded input by a tiny amount;
extra arithmetic precision alone cannot remove a frozen input defect.
The larger continuation residuals are not removed by this point test.


\subsection{Recursive finite-jet recalibration}
\label{subsec:v410-recursive-calibration}

The preceding frozen-data failures do not persist under all joint
changes of lower jets. A subsequent comparison holds the same field
values, original backgrounds and couplings fixed and first solves for
compatible first and symmetric second jets. Releasing second jets then
allows a symmetric third-jet completion while retaining those first
jets. The scalar condition
\(g=4-h\Psi^\dagger(Z\otimes\beta)\Psi/2=0\) explains a necessary
part of the hierarchy: the first fit enforces \(g,dg\), but its second
jets need not satisfy \(d^2g=0\). Free third jets cannot repair that
lower-data defect. The joint second/third fit does repair it numerically.

A reusable construction stores Taylor coefficients
\(c_\alpha=\partial^\alpha\Psi/\alpha!\), with multiindex factorial.
At each tested level the highest jets \(z_p\) enter linearly:
\[
 A_p z_p+b_p(J_{<p})=0,\qquad L_p A_p=0.
\]
A complete left kernel eliminates \(z_p\). The numerical search adjusts
allowed lower jets to satisfy \(L_p b_p=0\) together with every retained
lower-order residual, then recovers \(z_p\) and checks the full equations.
This is a construction algorithm with a possible failure outcome, not
an all-order existence assertion. Small updates and coefficient guards
do not prove convergence or uniqueness.

For the fourth-order test, field values and all first jets remain fixed;
second and third jets vary and fully symmetric fourth jets are recovered.
The highest matrix has size \(480\times840\), numerical rank \(424\),
and \(56\) left compatibility directions in the declared coefficient
packing. Five accepted steps at each of three points produce the raw
symmetric derivative norms below:
\begin{center}
\begin{tabular}{@{}lrrrr@{}}\toprule
point & \(\|F\|\) & \(\|dF\|\) & \(\|d^2F\|\) & \(\|d^3F\|\)\\
\midrule
centre & \(3.1\,10^{-14}\) & \(3.9\,10^{-14}\) & \(2.4\,10^{-13}\) & \(6.4\,10^{-12}\)\\
edge & \(2.0\,10^{-15}\) & \(2.9\,10^{-14}\) & \(2.5\,10^{-13}\) & \(5.4\,10^{-12}\)\\
inner & \(2.6\,10^{-15}\) & \(3.8\,10^{-14}\) & \(2.3\,10^{-13}\) & \(6.5\,10^{-12}\)\\
\bottomrule\end{tabular}
\end{center}
Normalized coefficient norms and raw derivative norms are stored
separately. A separate NumPy directional-series evaluator with a
recurrence inverse checks coefficients through degree three in six
directions per point, with discrepancies below \(4\,10^{-13}\).
Selected highest-operator columns and scalar constraint coefficients
are also checked. This is bounded implementation evidence, not a
certified exact root or an exhaustive theorem check.

Raw second-, third- and fourth-jet norms are approximately \(26,89,900\)
in the retained unscaled coordinates. They establish no physical scale
and no uniform growth estimate. The three germs are not shown to belong
to one regional solution; stabilization of lower coefficients across
levels, convergence, regional existence and native physical identification
remain open. The tested reusable construction and earlier failed controls
are archived under \texttt{analysis/op5-field-slices/}, in the dated
state-sufficiency directory. Neither these finite completions nor the
algorithm closes OP5.

\paragraph{Status of the 4.10 consolidation.}
Exact reconstruction and local operation identities are separated from
finite numerical evidence. Native action, clock identification, complete
octonionic structural embedding, regional spinor existence and OP5 closure
remain open. Scripts, inputs, coefficients, corrections and individual
scope statements are in \texttt{analysis/op5-field-slices/}.
The guard and global residual records are not promoted by these tests.
